


Date: 2026-11-15. Branch \texttt{round5} from tag \texttt{algo4-v1}. Key prefix \texttt{r5.}. This file is committed before any development run of any algorithm. No tagged result and no existing key is changed. Hardware data (\texttt{data/hardware}) is not used. The results of this round are reported as a separate round 5; they never change an earlier verdict. Session 1 covers G0 and L1 to L5; session 2 covers L6 to L10, R and D. For L6 to L10 this file fixes the hypotheses, the roles rules and the statistics; the choice of the datasets and the hashed splits is written as an addendum (V16 addendum S2) and committed before any session 2 run.

\subsection*{0. What was looked at before this file, and what this file corrects}

Looked at before this file: code and structure only (module headers, the column layout of the stored event table, the ns-3 application source, the sensor datasheet notes). No development or held-out number of any L algorithm had been computed. The baseline test suite of the repository (421 passed, 2 skipped) was run on the clean \texttt{algo4-v1} state.

Points where the task text and the repository disagree, written down now so that they are on record before any run:

\begin{itemize}
\item \textbf{L1, axis flips.} The earlier simulated variants \texttt{F\_x} and \texttt{F\_y} of \texttt{hwsim\_v2.py} change the sign of one axis of the acceleration and of the gyro. That is a mirror image of a right-handed sensor, not a rotation; no mounting of a real board does it. No method that uses gravity and the direction of travel can recover the left and right sides under a mirror image without an assumption about handedness. The physical readings of "a flipped axis" are 180-degree rotations (about z: the board is turned round, forward and sideways both reverse; about x or y: the board is upside down with one horizontal axis reversed). The primary held-out set of L1 uses those rotations. The mirror variants are run as a descriptive row and are not part of the criterion of H-L1.
\item \textbf{L1, calibration signal.} The earlier simulated event logs (3 s long, impact at about 1 s) hold no still period of 2 s and no straight-line braking before the impact. The calibration drive is therefore generated as a separate recording of the same installation (stillness, then a straight braking or acceleration, with GPS speed at 1 Hz). The NHTSA crash pulses are real; their rotations and the calibration drive are synthetic (stated in every sentence that uses them).
\item \textbf{L2, H-L2b.} With the provisional alert unchanged, a stricter confirm rule can lower the retraction share only by confirming provisional alerts that the current confirm rule rejects (crashes with a low final detector score but a strong post-event signature); it raises the share when it cancels provisional alerts that the current rule confirms. The hypothesis measures the net effect; both directions are reported.
\item \textbf{L2, time.} The corroboration window starts after the full window (1.8 s after the trigger) and lasts at most 3 s, so a confirmed alert leaves at most 3 s later than today. The provisional alert is unchanged. The cost is reported.
\item \textbf{L3, buffer and idle rate.} A buffer at the full rate and an idle sensor at a low rate cannot both hold, because the samples before the interrupt exist only at the idle rate. The design below is the reading that is physically consistent: in idle the accelerometer runs alone and fills the sensor FIFO (1,024 bytes = 170 accelerometer samples of 6 bytes, 1.7 s at 100 Hz), the gyro is in standby and the MCU is in light sleep; the motion interrupt wakes the MCU and turns the gyro on. The pre-trigger gyro window is therefore missing (zero) and the earliest samples after the wake-up lack the gyro start-up time. Reduced accelerometer idle rates (25 Hz and 10 Hz, zero-order hold to the 100 Hz grid) are run as sensitivity arms.
\item \textbf{L4, the link estimate.} ns-3 end devices in the application (\texttt{xsim\_alert.cc}) have no downlink SNR memory (heartbeats are unconfirmed). The start spreading factor is chosen from a link-margin estimate that equals the true received-power margin of the vehicle plus a Gaussian error of 2 dB; this is an assumption and the verdict says so. The \texttt{quick} policy of the repository already steps the spreading factor up after two failed rounds (gap 0.3 s); it is the nearest earlier method and is run as an additional baseline (descriptive), next to the zero-spacing control required by the task.
\item \textbf{L5, censoring.} About 60\% of the crash events alarm at the 1\% target, so the 90th percentile of the time to awareness over all crash events is the cap itself in both arms. The statistic is therefore defined with a cap and an additional conditional version (section 6).
\end{itemize}

\subsection*{1. G0: integrity rules (these are binding for every algorithm)}

\begin{itemize}
\item One Holm family for this round of at most 14 hypotheses (13 are used, section 9). It is reported separately; earlier verdicts never change.
\item Every algorithm has fixed data roles (section 2): DEVELOPMENT (tuning, at most 5 logged iterations per algorithm), then the sources and the parameter file are hashed (freeze), then the HELD-OUT run once. A held-out set of the round was not held-out in rounds 3 or 4; where a whole set was (the VZCrash test split, scored by \texttt{s6\_early} and \texttt{r3\_t3}), a disjoint half by hashed event id is used (section 2.1). For new datasets the split is by subject, vehicle, device or year, and the hash of the id lists is written before the held-out run.
\item Code verification for every new module before the freeze (section 10): unit tests; property tests; a known-answer reproduction of one earlier number; at least 20 mutants with at least 80\% killed; a review by a fresh-context reviewer against this file.
\item An algorithm counts as a contribution only if its primary hypothesis is supported on held-out data (Holm-adjusted) AND its evidence label is MEASURED\_SOFTWARE, TRACE\_DRIVEN or FORMAL, or it holds in at least two link models. Results that are SIMULATED only or SIMULATED\_HW\_FORMAT only go to the supplement.
\item Every algorithm gets a row in \texttt{docs/06\_references/NOVELTY\_MATRIX\_2021PLUS.md} before it is run (OpenAlex, journal articles from 2021, plus older and conference context; the closest items opened; overlap level). If the overlap is substantial, the closest prior method becomes a baseline.
\item Words: the words banned by \texttt{CLAUDE.md} are not used in generated text. A built-but-unmeasured item is called "built, not yet measured".
\end{itemize}

\subsection*{2. Data roles}

\subsubsection*{2.1 VZCrash (CC BY-NC 4.0) and the hashed halves}

The de-duplicated splits of the repository are used as they are (train, validation, test; \texttt{windows\_cache}, \texttt{split.py}). The test split was scored by \texttt{s6\_early} and used again by \texttt{r3\_t3}, so it is not held-out for round 5 as a whole. The test split is divided by the hash rule \texttt{half(id) = int(sha256("r5|" + event\_id), 16) mod 2} (event id as decimal text). \textbf{HALF1} is the held-out set of L2; \textbf{HALF0} is the held-out set of L3 and of L5. The other half of each algorithm is not touched by that algorithm. Train and validation are DEVELOPMENT for every algorithm (train: fitting; validation: thresholds, selection of the iteration). Each algorithm touches its half once, after its freeze. The hashes of the sorted event-id lists of both halves are written to \texttt{results/tables/r5\_split\_hashes.json} before the initial held-out run.

\subsubsection*{2.2 L1 data}

{\scriptsize\raggedright
\begin{longtable}{p{0.475\linewidth}p{0.475\linewidth}}
\toprule
Role & Content \\
\midrule
\endhead
DEVELOPMENT, simulated logs & \texttt{hwsim\_v2} generators with the base seed 50261115 (installations D01 to D12, section 3.2); never the seeds of V11 (20261114 + 1000 i) \\
HELD-OUT, simulated logs & the same generators with the base seed 60261116 (installations H01 to H12); regenerated, hashed, run once \\
DEVELOPMENT, NHTSA & the groups with role \texttt{tuning} of the make-and-model split A of \texttt{direction.py} (\texttt{split\_groups}, as in X1) \\
HELD-OUT, NHTSA & the groups with role \texttt{heldout} of the same split; real pulses, synthetic rotations drawn with seed 42 (section 3.3) \\
\bottomrule
\end{longtable}}

\subsubsection*{2.3 L3 data}

DEVELOPMENT: VZCrash train and validation (the detector is the frozen forest family of R3, retrained per seed exactly as \texttt{s6\_early} does). HELD-OUT: HALF0 of the test split (detection), the NHTSA \texttt{heldout} groups (side accuracy, as in L1 and with the same rotations disabled), and the normal-driving events of the VZCrash validation split for the measured active share used by the energy estimate (a measured input to an analytical formula).

\subsubsection*{2.4 L4 data (SIMULATED)}

ns-3 application of the repository (V13 settings: 60 vehicles, 10 events per vehicle-hour, duty 1\%, ADR and application backoff off, one gateway at the centre). DEVELOPMENT seeds 401 to 420, HELD-OUT seeds 501 to 540 (fresh: not used before). Grid: radius 1,000, 1,500, 2,000, 2,500 and 3,000 m; the two calibration sensitivity settings (\texttt{calib}, \texttt{shipped}); extra loss per direction 0, 0.1 and 0.3 (30 cells); the base spreading factor of the control is SF7. The custom simulator uses the same grid in its own units.

\subsubsection*{2.5 L5 data}

Events and detector scores: VZCrash train (fit) and validation (choice of the operating point) as DEVELOPMENT, HALF0 of the test split as HELD-OUT. Network delay distributions: custom simulator (SIMULATED), ns-3 (SIMULATED; fresh seeds 601 to 640 for the held-out evaluation, 421 to 440 for development), LoED traces (TRACE\_DRIVEN: round-3 \texttt{half = 0} devices for development, \texttt{half = 1} devices for the held-out evaluation; the delay is an input to the evaluation, not the quantity tested).

\subsubsection*{2.6 L6 to L10}

Roles are fixed in V16 addendum S2 before any session 2 run, by these rules: L6 fits on earlier years and tests on the latest year of a public geocoded crash dataset; L7 and L8 have no data split (FORMAL and enumeration; L8b replays held-out VZCrash event rates, HALF0 of the test split); L9 uses stored results (no new held-out set); L10 uses published splits of the benchmarks audited, with subject ids where they exist.

\subsection*{3. L1: gravity-referenced automatic mounting calibration (SIMULATED\_HW\_FORMAT and MEASURED\_SOFTWARE)}

\subsubsection*{3.1 Design (frozen here; development may change only the items in 3.5)}

Input: the logged accelerometer and gyro samples (100 Hz) of a calibration drive of the same installation, GPS speed (1 Hz) and the crash log.

\begin{itemize}
\item \textbf{Stillness.} A window of 2.0 s is still if the gyro magnitude is below 3 deg/s at every sample, ||a| - 1 g| is below 0.05 g at every sample and the standard deviation of |a| is below 0.02 g. The gravity direction \texttt{z\_hat} is the unit vector of the mean acceleration over the latest still window (specific force at rest points up).
\item \textbf{Forward axis.} A segment of at least 2.0 s in which the GPS speed changes by at least \texttt{A\_MIN} = 1.0 m/s per second on average (monotonically at the 1 Hz samples), the gyro component along \texttt{z\_hat} stays below 5 deg/s (straight line) and the horizontal part of the mean acceleration (the component of the mean acceleration perpendicular to \texttt{z\_hat}) is at least 0.05 g. The horizontal direction \texttt{d} of that mean acceleration is the forward axis when the speed rises and the backward axis when it falls: \texttt{x\_hat = sign(dv/dt) * d / |d|}.
\item \textbf{Vehicle frame.} \texttt{y\_hat = z\_hat x x\_hat} (cross product; this is the sign convention of the native axes of the repository, in which the stored calibration signs +1 and -1 apply). The rotation \texttt{R\_hat} has the rows \texttt{x\_hat}, \texttt{y\_hat}, \texttt{z\_hat}; the sample in the vehicle frame is \texttt{R\_hat s}.
\item \textbf{Side rule.} The existing rule (\texttt{sh\_events.side\_rule}, impulse over 60 ms after the 2 g trigger, calibration signs +1 and -1) is applied to the rotated samples. The rule is unchanged.
\item \textbf{Re-estimation.} At every still window the gravity direction is compared with the stored \texttt{z\_hat}; if the angle exceeds 10 degrees the calibration is marked stale and the side rule withholds a decision (reported as "no decision") until a new forward estimate exists. A decision with a stale calibration counts as wrong in the accuracy figures.
\item Comparison (control): the unchanged rule on the raw sensor axes.
\end{itemize}

\subsubsection*{3.2 Simulated logs (SIMULATED\_HW\_FORMAT)}

The generators of \texttt{hwsim\_v2.py} are used through a new module (\texttt{r5\_l1.py}) that adds a general rotation of the mounting and a calibration drive; with the identity rotation it reproduces the sensor rows of \texttt{hwsim\_v2.sense} for the same random stream (known-answer test). Installations (12 per role): identity control; tilt about x of 15 and 30 degrees; tilt about y of 15 and 30 degrees (as the earlier variants); rotation by 180 degrees about z, about x, about y; rotation about z by +90 and -90 degrees; two random mountings (a tilt of at most 30 degrees about a random horizontal axis combined with a yaw drawn from 90, 180 and 270 degrees); one combined 180 degrees about x with a 20-degree tilt. Each installation has 80 horizontal crash events (20 each for front, rear, left, right) and one calibration drive (3 s still, then braking at 0.25 to 0.45 g for 3 s with the GPS speed at 1 Hz and 0.1 m/s noise, or an acceleration of the same size in half of the installations). Peak, noise, bias, filter and rounding are those of \texttt{hwsim\_v2}. Mirror variants (\texttt{F\_x}, \texttt{F\_y} semantics) are run as an additional descriptive row.

\subsubsection*{3.3 NHTSA part (MEASURED\_SOFTWARE for the pulses, synthetic rotation and calibration drive)}

Held-out tests (frontal barrier, side impactor, oblique) of the \texttt{heldout} groups, 16 g range, 100 Hz sensor model, 4 random rotations per test (seed 42; yaw from 0, 90, 180, 270 degrees, a tilt of at most 30 degrees about a random horizontal axis with magnitude uniform on 0 to 30, and an additional 180 degrees about x with probability 0.25). The rotation is applied to the raw record before the sensor model (so the clipping acts per sensor axis). The calibration drive of the same rotation is synthetic (as in 3.2). Score: balanced accuracy, front against side, sign-free, as X1 (\texttt{direction.balanced\_acc}); the left-against-right sign is not scored (the NHTSA side tests all strike the left).

\subsubsection*{3.4 Hypothesis H-L1}

H-L1: with L1 the side accuracy is at least 0.90 under every held-out physical installation, and the NHTSA balanced accuracy is at least 0.90. Statistic for the Holm family: the one-sided exact binomial test of H0 "pooled side accuracy at most 0.90" over the pooled crash events of the 13 held-out simulated installations (physical set, mirror rows excluded). Criterion for "criterion met": the point accuracy of each of the 12 installations is at least 0.90 AND the lower 95\% limit of the pooled accuracy (cluster bootstrap over events, 2,000 resamples, seed 42) is at least 0.90. The NHTSA part passes if the pooled balanced accuracy is at least 0.90 and its lower 95\% limit (bootstrap over groups, 2,000 resamples, seed 42) is at least 0.85. \textbf{Per G0 rule 4, L1 is a contribution only if the NHTSA part also passes; otherwise it goes to the supplement.} The control (raw axes) is reported on the same installations.

\subsubsection*{3.5 Development (at most 5 logged iterations)}

Items development may change: \texttt{A\_MIN} (0.5, 1.0, 1.5 m/s per second), the still-window length (1.5, 2.0 s), the stillness gyro limit (2, 3, 5 deg/s), the stale-angle limit (5, 10, 15 degrees). Selection rule: the highest pooled development side accuracy; ties go to the iteration nearer the defaults of 3.1. If no iteration beats the defaults by more than 1 point, the defaults are kept.

\subsection*{4. L2: post-event state corroboration (MEASURED\_SOFTWARE)}

\subsubsection*{4.1 Design}

After the detector fires, the node waits a window of \texttt{W} seconds after the full window (decision at 1.8 s after the trigger) and computes post-event features from the raw event: the speed drop (GPS speed over the 3 s before the trigger, median, minus the speed at the sample nearest the end of the window; GPS samples are placed by proportional index over the event: sample j of k at j times (n/100)/k seconds, stated assumption); stillness (mean gyro magnitude in the window; the share of samples with ||a| - 1 g| below 0.05 g in the last second of the window); orientation change (the angle between the mean acceleration direction of the 1 s before the trigger and of the last 0.5 s of the window); rotation energy (sum of the squared gyro magnitude from the trigger to the end of the window). A small model (logistic regression on standardised features, or a tree of depth 3) combines the 1,800 ms detector score \texttt{s\_F} of the frozen forest (seed 42) with these features. Missing GPS leaves the speed features at the training median and sets an indicator. Trained on triggered train events (crash against near-miss and normal driving), thresholds on validation only. The provisional alert and the early-exit stage are unchanged; L2 only replaces the confirm or cancel decision.

\subsubsection*{4.2 Hypotheses (Holm family)}

Baseline: the confirm stage of R3 (full window alarm \texttt{F} at the validation threshold for the 1\% target), recomputed here by the same code as \texttt{s6\_early} (known-answer check against \texttt{r3.t3.a1.*}).

\begin{itemize}
\item \textbf{H-L2a.} At the crash recall of the baseline confirm stage (L2 threshold: the largest validation threshold whose validation crash recall, system level, is at least the baseline's), the near-miss false-alarm rate (alarms among all near-miss events of HALF1; events that do not trigger count as no alarm) falls by at least 25\% relative. Criterion: point reduction at least 0.25; lower 95\% limit of the reduction above 0 (paired bootstrap over events, 2,000 resamples, seed 42); and held-out crash recall of L2 at least the baseline's minus 0.01. One-sided p: share of resamples with reduction at most 0.
\item \textbf{H-L2b.} At the false-alarm budget of the baseline (L2 threshold chosen on validation so that the system false-alarm rate over all negatives equals the baseline's validation rate), the retraction share of provisional alerts (provisional alarms not confirmed, over provisional alarms) falls by at least 3 points (0.03) and the final crash recall falls by at most 1 point (0.01). Criterion: point fall at least 0.03, lower 95\% limit above 0 (paired bootstrap over events), and the recall condition at the point level. One-sided p: share of resamples with a fall at most 0.
\end{itemize}

Kill: both hypotheses fail. The extra confirm delay (the window \texttt{W}) is reported with the result.

\subsubsection*{4.3 Development (at most 5 logged iterations)}

Items: \texttt{W} in \{1, 2, 3\} s; the model (logistic, tree); the class weight (none, balanced). Selection rule: the iteration with the lowest validation near-miss false-alarm rate at the baseline's validation recall; ties go to the smaller \texttt{W}. The defaults are \texttt{W} = 2 s, logistic, no weight.

\subsection*{5. L3: wake-on-motion with a pre-trigger buffer (MEASURED\_SOFTWARE for detection, ANALYTICAL for energy)}

\subsubsection*{5.1 Design}

As corrected in section 0 point 5. Arms: A = always-on, 100 Hz accelerometer and gyro (the repository's pipeline). B100 (primary) = accelerometer-only FIFO ring at 100 Hz in idle (capacity 170 samples = 1.7 s; the detector needs 1.0 s), gyro in standby, MCU in light sleep; the motion interrupt fires at the earliest sample where the high-passed deviation of the acceleration magnitude from its running mean exceeds \texttt{T\_MOT} (0.5 g; the high-pass is the running mean of the previous 0.5 s); the gyro delivers valid samples \texttt{L\_GYRO} = 35 ms (4 samples at 100 Hz) after the interrupt (value from the MPU6050 datasheet, rechecked with the opened datasheet when the energy table is written); gyro samples before that are zero. Sensitivity arms: B25 and B10 (accelerometer idle at 25 and 10 Hz, zero-order hold to 100 Hz, otherwise as B100). The detector (forest, threshold) is the frozen R3 family retrained per seed on the windows of each arm (train), with the threshold from validation for the 1\% target.

\subsubsection*{5.2 Energy (ANALYTICAL)}

\texttt{E\_B / E\_A = (P\_idle (1 - d) + P\_act d) / P\_act}, where \texttt{P} is the sum of the sensor and MCU currents at the supply voltage, \texttt{d} is the share of time the interrupt is asserted plus a hold of 2 s, measured on the normal-driving events of the VZCrash validation split (16 s windows selected by the dataset's labelling, so \texttt{d} is likely high; the limit is stated) with a bootstrap interval over events, and the currents are the typical values of the opened MPU6050 and ESP32 datasheets with page numbers (\texttt{data/raw/t4\_datasheets} holds the ESP32 file; the MPU6050 note is \texttt{SENSOR\_DATASHEETS.md}). Where a current is not in an opened datasheet the energy result is NOT\_RUN for that arm.

\subsubsection*{5.3 Hypothesis H-L3 (Holm family)}

H-L3: on HALF0 the system crash recall of B100 at the 1\% target is within 1 point of A (recall of A minus recall of B100 at most 0.01), and the estimated sensor and MCU energy of B100 is at least 50\% below A (ratio at most 0.50, also at the upper 95\% limit of \texttt{d}). One-sided p: share of paired-bootstrap resamples (over events, 2,000, seed 42) with a recall loss of at least 0.01. If the energy part is NOT\_RUN the hypothesis is not supported. Side accuracy (NHTSA held-out groups, the stitched pulse at B100) and false-alarm rates are reported.

\subsubsection*{5.4 Development (at most 5 logged iterations)}

Items: \texttt{T\_MOT} (0.3, 0.5, 1.0 g); \texttt{L\_GYRO} (35, 70 ms); retrain the detector on the arm's windows or score them with the frozen forest of A. Selection: the lowest validation recall loss against A; ties go to the lower \texttt{d}.

\subsection*{6. L4: deadline-optimal spreading-factor ladder (SIMULATED and ANALYTICAL)}

\subsubsection*{6.1 Design}

The alert is the 97-byte confirmed alert; the deadline is 10 s from the crash; the vehicle sends again at once after each failed round (gap 0, as the zero-spacing control); the duty limit is 1\% (sliding hour). The ladder starts at the lowest SF whose margin (estimated received power minus the SF sensitivity) is at least \texttt{M0} dB, and steps up one SF (to SF12) after \texttt{k} consecutive failed rounds. The plan (start SF, \texttt{k}) is computed by an analytical routine that maximises the expected delivery within the deadline: attempts at SF \texttt{s} take the airtime plus the receive-window time (0.32 s); the per-attempt success of an SF is \texttt{(1 - p\_s) q\_s} with \texttt{p\_s = p 0.7\textasciicircum{}(s - SF\_base)} and \texttt{q\_s} from the margin. The estimate error is 2 dB (section 0 point 6). Controls: Z0 (fixed SF7, gap 0, same retries) and \texttt{quick} (descriptive).

\subsubsection*{6.2 Hypothesis H-L4 (Holm family)}

The "edge cells" are the cells of the grid (section 2.4) in which the Z0 control delivers within 10 s between 0.30 and 0.95 on the DEVELOPMENT seeds (all 30 cells are run on the development seeds with Z0 for this purpose). On the HELD-OUT seeds, paired by seed: the mean difference (ladder minus Z0) in delivery within 10 s over the edge cells is at least 0.03 with a lower 95\% limit above 0 (cluster bootstrap over runs, 2,000 resamples, seed 42), AND in no cell of the 30 the lower 95\% limit of the difference is below -0.01. One-sided p: share of resamples with a mean difference at most 0. The hypothesis must hold in both simulators (ns-3 and the custom simulator); if the custom-simulator part is not run it is NOT\_RUN and the hypothesis is reported as ns-3 only (supplement, G0 rule 4). Real traces cannot change the spreading factor, so no trace result is claimed.

\subsubsection*{6.3 Development (at most 5 logged iterations)}

Items: \texttt{k} in \{1, 2, 3\}; \texttt{M0} in \{0, 3, 6\} dB. Selection: the highest mean development delivery over the edge cells; ties go to the smaller \texttt{k}.

\subsection*{7. L5: end-to-end deadline-aware threshold co-design (MEASURED\_SOFTWARE plus the link-model label)}

\subsubsection*{7.1 Design}

Time to awareness of a crash event: the arrival at the base of the earliest frame that reaches it. Provisional decision at the checkpoint time of the early-exit rule (50 to 800 ms after the trigger) plus the delay of the 14-byte flash; confirm decision at 1,808 ms plus the delay of the 97-byte alert; for an event that does not alarm, that is not delivered within 10 s of the decision or that does not trigger, the time is the cap \texttt{T\_CAP} = 60 s (the responders learn by another route). Delays are drawn from the empirical delay distributions of the three link models at the primary scenario of each (SF9, extra loss 0.1 per direction; ns-3 \texttt{shipped} setting; LoED devices of the role's half), independently per event, 20 draws per event, seed 42. The co-designed operating point is the member of the grid of \texttt{early.py} (\texttt{A\_HI} x \texttt{B\_LO}, plus no early exit, with the final threshold \texttt{t\_f} the smallest that keeps the validation system false-alarm rate within the budget) that minimises the validation 90th percentile of the time to awareness subject to the same false-alarm budget as the baseline (1\% target: 0.10 false alarms per vehicle-hour at 10 events per vehicle-hour; the 5\% target is a sensitivity run). The baseline is the frozen R3 rule of \texttt{early.choose\_rule} at the same budget. The duty limit does not bind at these false-alarm rates and is checked, not optimised.

\subsubsection*{7.2 Hypothesis H-L5 (Holm family)}

H-L5: on HALF0 events and held-out delays, the co-designed operating point lowers the 90th percentile of the time to awareness (with the cap) by at least 10\% relative to the baseline at the 1\% budget, in at least two of the three link models. Per link model: point reduction at least 0.10 and a lower 95\% limit above 0 (paired bootstrap over events, 2,000 resamples, seed 42). If the cap is the 90th percentile in both arms (recall below 90\%) the cell is "tied at the cap" and the conditional statistic is used instead: the 90th percentile among the events that were alerted and delivered. The hypothesis p-value is the second smallest of the three link-model one-sided p-values (a two-of-three rule; the bootstrap share of resamples with a reduction at most 0), reported with the three values. The label of each link model is stated (custom SIMULATED; ns-3 SIMULATED; LoED TRACE\_DRIVEN).

\subsubsection*{7.3 Development (at most 5 logged iterations)}

Items: the grid subset searched (full grid, \texttt{A\_HI} only, \texttt{B\_LO} only); the final-threshold rule (smallest \texttt{t\_f}, or the \texttt{t\_f} that minimises the P90); whether the objective is the P90 with the cap or the conditional P90. Selection: the lowest validation objective; ties go to the smaller grid.

\subsection*{8. L6 to L10 (session 2; hypotheses fixed now)}

\begin{itemize}
\item \textbf{L6.} H-L6: with the same k, crash-weighted placement of receivers covers at least 10 percentage points more crashes of the held-out (latest) year than population-weighted placement, for every k in the pre-registered range (k = 3, 5, 8, 12; the range is fixed here). Comparisons: uniform grid, population-weighted, road-length-weighted. Path loss: log-distance with the exponent and shadowing fitted on LoED reception data (stated). MEASURED\_SOFTWARE for the crash data, ANALYTICAL for coverage.
\item \textbf{L7.} H-L7a (FORMAL): Tamarin or ProVerif proves that compromise of the current ratchet key reveals no earlier alert key and allows no forgery of an earlier alert (forward secrecy), and that authenticity and the per-vehicle ACK binding still hold. H-L7b (exhaustive enumeration): resynchronisation succeeds after up to W lost ACKs inside the look-ahead window (W fixed in the addendum). Cost in bytes and hash operations (ANALYTICAL).
\item \textbf{L8.} H-L8a (FORMAL or bounded model): a single compromised vehicle cannot cause more than R alerts per hour to be forwarded as high-confidence. H-L8b (SIMULATED plus real VZCrash event rates): no legitimate crash alert is quarantined in the held-out replay. Supplement unless H-L8a is FORMAL.
\item \textbf{L9.} H-L9 (descriptive, not in the Holm family): the break-even cost ratio between a false dispatch and a minute of delay, with intervals, from the stored T3 retraction rates and time gains and from L2 if it passes. No superiority claim.
\item \textbf{L10.} H-L10a: at least one more public benchmark has more than 1\% of its test windows identical (exact hash and the window-identity rule of C1, unchanged) to a window in another split. H-L10b: where a published baseline can be reproduced, removing the leaked windows changes its headline metric by more than the published interval or 1 point. If no other benchmark shows leakage, that is reported as a finding that narrows C1.
\end{itemize}

\subsection*{9. Multiple testing and decision rule}

One Holm family of 13 hypotheses: H-L1, H-L2a, H-L2b, H-L3, H-L4, H-L5 (session 1) and H-L6, H-L7a, H-L7b, H-L8a, H-L8b, H-L10a, H-L10b (session 2); H-L9 is descriptive. In session 1 the adjusted p-values use \texttt{m = 13} with the session 2 hypotheses entered as p = 1 (which keeps the step-down valid and conservative); the final adjusted values are computed at the end of session 2. A hypothesis is supported if its criterion holds at the point level and the Holm-adjusted p is below 0.05. FORMAL hypotheses are decided by the proof outcome and enter Holm with p = 0 (supported) or p = 1 (not supported or undecided). A kill rule that fires is applied whatever the other results are.

\subsection*{10. Code verification (report \texttt{docs/03\_results/R5\_CODE\_VERIFICATION.md})}

For each new module: unit tests; property tests (L1: the estimated rotation is a proper rotation; the corrected side result is invariant under a rotation applied to calibration and events together when noise is zero; L2: no feature uses samples before the trigger window start except the stated pre-trigger ones, and none after the end of the window; L3: the stitched window never uses a sample after the interrupt for the pre-trigger part and the FIFO capacity bound holds; L4: the ladder never plans an attempt that ends after the deadline or beyond the duty limit, and the start SF is monotone in the margin; L5: the chosen point meets the false-alarm budget on validation); a known-answer reproduction of an earlier number (L1: \texttt{hwsim\_v2.sense} rows for the identity rotation, and the control side accuracy on V11 \texttt{V00\_base} and \texttt{F\_y}-free variants against \texttt{sh.*} keys; L2 and L3: the baseline recall and the retraction share against \texttt{r3.t3.a1.*}; L4: the ns-3 policy \texttt{z0} against the stored \texttt{quick}/\texttt{p1} mechanics on three stored seeds, plus the analytical attempt count against the ns-3 invariant; L5: the baseline arm against \texttt{early.choose\_rule} output); at least 20 mutants per algorithm (hand-written, as \texttt{scripts/a4\_mutation.py}) with at least 80\% killed; a review of the code against this file by a reviewer with no access to the results, before the freeze. The ns-3 application change is checked by byte-identical outputs of the earlier policies on three stored seeds and by the per-run invariants of the new policies (no transmission ending after the deadline; hourly airtime never above the limit).

\subsection*{11. Literature rows}

For L1 to L5 the rows of \texttt{NOVELTY\_MATRIX\_2021PLUS.md} are written before the algorithm's held-out run, from OpenAlex queries (journal articles from 2021-01-01, filter on source type), with older and conference items as context. An absence is written only as "no journal article from 2021 onwards found in the sources checked (OpenAlex; IEEE Xplore returns HTTP 418 to scripts and is not run; Scopus, Google Scholar and TRID not run)". Closest items are opened at abstract level, and full text where open access allows.

\subsection*{12. Seeds, code and files}

Seed 42 for every draw; seeds 0 to 4 for repeated model fits (the held-out comparison of L2, L3 and L5 uses seed 42 as the primary and reports seeds 0 to 4 as a sensitivity range). Code: \texttt{r5\_common.py}, \texttt{r5\_l1.py}, \texttt{r5\_l2.py}, \texttt{r5\_l3.py}, \texttt{r5\_l4.py}, \texttt{r5\_l5.py}, \texttt{r5\_decide.py}, the ns-3 policies \texttt{z0} and \texttt{l4} in \texttt{xsim/ns3/xsim\_alert.cc}. Results under \texttt{r5.} in \texttt{results/results.json} and \texttt{results/tables/r5\_*}. Frozen files and parameter hashes: \texttt{results/tables/r5\_frozen.json}.

\subsection*{13. Addendum A1 (committed before the initial development run of any algorithm)}

Two corrections found while the code was being planned, with the MPU-6050 datasheet (Rev 3.3, 2012, opened 2026-11-15 from the copy named in \texttt{SENSOR\_DATASHEETS.md}) read for the two values; no development or held-out number of any algorithm existed at the time.

\begin{itemize}
\item \textbf{L1 stillness gate.} Section 3.1 item 1 sets the gate on ||a| - 1 g| at 0.05 g. The datasheet gives a zero-g initial tolerance of +-50 mg on x and y and +-80 mg on z (page 13), and the simulated installations use a bias of 0.03 g per axis (standard deviation), so a gate of 0.05 g would reject about one installation in ten for the bias alone. The gate is set to \textbf{0.15 g} (above the datasheet tolerance with room for noise). The other limits of 3.1 are unchanged. The gate is not a development item.
\item \textbf{L3 gyro start-up.} Section 5.1 gives \texttt{L\_GYRO} = 35 ms. The datasheet gives a gyroscope start-up (zero-rate settling to +-1 deg/s from power-on) of \textbf{30 ms} (page 12). The default is 30 ms (3 samples at 100 Hz); the development alternative stays 70 ms. The same page gives: gyroscope only 3.6 mA, accelerometer only 500 uA, gyroscope and accelerometer 3.8 mA (page 14); the FIFO holds 1,024 bytes (page 7).
\item \textbf{L1 installation count.} Section 3.2 lists 13 installations (identity control; four tilts; three 180-degree rotations; two yaws; two random mountings; one combined), not 12. The criterion of H-L1 applies to every one of the 13 (the identity control included), and the pooled figures use the 13.
\end{itemize}


\subsection*{14. Addendum A2 (after the independent code review, before the initial development run of any algorithm)}

A reviewer with no access to results read the code of L1 to L5 against this file. The review found points that the code now follows and points where this file is made precise. No development or held-out number existed at the time.

\begin{itemize}
\item \textbf{L3 awake state (blocker).} The initial draft treated the node as awake from the earliest interrupt of the whole 16 s clip, which contradicts the 2 s hold of the energy model. The node is awake from the start of the awake period that contains the 2 g trigger sample: an assertion keeps it awake for 2 s (200 samples); an earlier bump more than 2 s before the event does not. If the node is idle at the trigger sample the event is missed by B. The detector of B is retrained on the arm's windows by default (section 5.1); scoring with the frozen forest of A is development iteration 4. The sensitivity arms B25 and B10 detect the interrupt on the 100 Hz signal (they assume that a pulse of 10 to 45 ms still wakes an accelerometer sampling at 10 to 25 Hz, which is optimistic) and are labelled so.
\item \textbf{L3 not run in this session:} the side accuracy of the stitched NHTSA pulse is equal to the always-on value by construction (the accelerometer samples of B100 are those of A) and is not computed (NOT\_RUN, reason stated in the result key).
\item \textbf{L5.} (a) The final alert is the confirm alarm of the 1,800 ms forest at its validation threshold (R3); the "no early exit" candidate and the final-threshold rule \texttt{t\_f} of section 7.3 have no effect on provisional alerts and are dropped; the development items are the grid subset and the objective. (b) A frame must reach the base within 10 s of the crash (the R3 convention). A provisional alarm that the confirm stage rejects is cancelled and the event counts as not aware even when the flash was delivered. (c) The conditional statistic is used in the held-out test only when both arms are at the cap; the objective item steers development only. (d) The development ns-3 runs use an event table built from VALIDATION decisions; the held-out runs use the table of the test split (\texttt{pp\_r3\_ns3.make\_events}). (e) The development selection compares the mean validation reduction of the iterations; an iteration with the conditional objective is compared on its own objective, which gives a reduction even when the cap statistic is tied, so that iteration can win by construction; the result states which iteration was chosen. (f) The held-out delay draws use seed 42.
\item \textbf{L1.} (a) Re-estimation (section 3.1 item 5) is implemented and unit-tested only; the 3 s crash logs contain no still segment, so no installation produces a stale gravity direction and the development item \texttt{ang\_stale} has no effect on any result. (b) The pooled lower limit of the accuracy resamples installations (clusters) with replacement, because the events of one installation share one calibration. (c) The gyro straight-line test uses the segment mean, and the speed condition requires every 1 Hz interval to reach \texttt{A\_MIN}; the road vibration of the synthetic drive is 0.003 g (standard deviation), not 0.008 g. These are implementation readings of "on average".
\item \textbf{L4.} (a) The custom-simulator part is NOT\_RUN in this session (the cut order allows it): H-L4 is reported as ns-3 only. (b) Events dropped because the vehicle is still active are removed from the delivery share of the run (the sets can differ slightly between Z0 and L4). (c) The invariants of the new policies are counted at the PHY (end of a transmission after the deadline; trailing-hour airtime above the limit, from the independent \texttt{DUTY} line). (d) The analytic start-SF rule does not use the shipped sensitivity offsets; it is used in the analytical part only.
\item \textbf{Seeds 0 to 4 (section 12)} and the 5\% budget run of L5 are NOT\_RUN in this session; the primary results use seed 42.
\item \textbf{Freeze.} The freeze file records the sources (including \texttt{r5\_common.py} and the older modules each algorithm calls), the cached data files, the parameters and the git commit; a held-out run refuses to start if its held-out keys already exist; the hashes of the two test halves are written to \texttt{results/tables/r5\_split\_hashes.json} before the initial held-out run.
\end{itemize}


\subsection*{15. Addendum S2 (session 2: datasets, roles, splits, hashes, designs of L6 to L11; committed before any session-2 run)}

Session 1 is closed (branch \texttt{round5}, commit c533077); nothing from it is rerun or re-tuned. The round-5 Holm family stays at the 13 of section 9 (the task text of session 2 says 14; the count used in session 1 is kept and this is stated in the paper). L11 was added after session 1 because session 1 showed that L1 does not recover mirror-image axis flips. \textbf{L11 is labelled ADDED-AFTER-SESSION-1 in every table and is reported in its own one-hypothesis family (H-L11a), outside the family of 13.} H-L11b is secondary and simulated (supplement).

\subsubsection*{15.1 Downloaded data (hashes of the files as downloaded on 2026-11-15; raw data is git-ignored)}

{\scriptsize\raggedright
\begin{longtable}{p{0.237\linewidth}p{0.237\linewidth}p{0.237\linewidth}p{0.237\linewidth}}
\toprule
Use & File & SHA-256 (start) & Licence \\
\midrule
\endhead
L6 crashes & FARS 2019 to 2022 national CSV zips (\texttt{static.nhtsa.gov}), file \texttt{accident.csv} & \texttt{0974abb9}, \texttt{b2806902}, \texttt{743c19a1}, \texttt{989448d7} & US government work, public \\
L6 population proxy & GeoNames \texttt{cities15000.zip} (places of at least 15,000 inhabitants) & \texttt{ce7a4055} & CC BY 4.0 \\
L10 benchmark 1 & UCI Human Activity Recognition Using Smartphones (\texttt{uci\_har.zip}) & \texttt{c00b8030} & CC BY 4.0 \\
L10 benchmark 2 & SisFall (Zenodo 22212281, \texttt{SisFall.zip}) & \texttt{dd0f9cf1} & CC BY 4.0 \\
\bottomrule
\end{longtable}}

\subsubsection*{15.2 L6: crash-weighted gateway placement}

Crash data: FARS \texttt{accident.csv} with valid coordinates (latitude 24 to 50, longitude -125 to -66, the contiguous US). \textbf{DEVELOPMENT years 2019, 2020, 2021; HELD-OUT year 2022} (split by year, fixed here). Weights: number of fatalities (\texttt{FATALS}) of the crash, a severity weight. Candidate sites: the distinct 0.02-degree cells (about 2 km) holding at least one development crash, plus the GeoNames places; no 2022 information is used to build candidates. Methods: (a) crash-weighted greedy maximum coverage of the development crashes (fatality weights); (b) population-weighted greedy maximum coverage of the GeoNames places (weights: population; the proxy is places of at least 15,000 inhabitants, stated); (c) uniform lattice (the candidate sites nearest to a regular 1-degree lattice, the k that come earliest in a fixed order). The road-length-weighted comparison is NOT\_RUN (no road-length data was obtained). \textbf{Coverage:} a crash is covered if it lies within \texttt{R\_COV} = 4.0 km of a gateway; \texttt{R\_COV} follows from the log-distance model of the ns-3 application (14 dBm, 10 dB vehicle loss, 25 + 30 log10 d dB, no shadowing) and the SF9 sensitivity of -129 dBm (ANALYTICAL assumption). \textbf{The calibration on LoED reception data named in the task is NOT\_RUN:} the LoED records carry no device or gateway positions, so no distance can be tied to a reception. Sensitivity runs at 2.5 km (SF7) and 6.0 km. \textbf{k in \{50, 100, 200, 400\}.} H-L6: for every k, held-out (2022) coverage of method (a) exceeds that of (b) by at least 0.10 (absolute); statistic per k: paired bootstrap over 2022 crashes (2,000 resamples, seed 42) of the coverage difference; criterion: point at least 0.10 and lower 95\% limit above 0 for every k; p = the largest one-sided bootstrap p over the four k (intersection test). No tuned parameter (greedy is deterministic); the development years are used only to place the gateways.

\subsubsection*{15.3 L7: forward-secure ratchet (FORMAL)}

Design: per-vehicle alert key \texttt{K\_0} derived from the master key and the vehicle id; after each acknowledged alert and at each boot epoch \texttt{K\_\{i+1\} = H(K\_i, "ratchet")}; the alert tag is \texttt{MAC(K\_i, alert)}; the base keeps \texttt{K\_i} and derives at most \texttt{W} = 8 steps ahead; old keys are erased on both sides. \textbf{H-L7a} (Tamarin 1.12.0 in WSL, time box 3 hours in total): lemma 1 forward secrecy: if the base accepts an alert with index \texttt{i} of vehicle \texttt{v} that \texttt{v} did not send, then the attacker had compromised \texttt{v} at an index at most \texttt{i}; lemma 2 authenticity of accepted alerts without compromise; lemma 3 per-vehicle ACK binding (an ACK of one vehicle is not accepted by another); lemma 4 executability. Supported only if lemmas 1 to 3 are proved; "undecided" counts as not supported (p = 1). \textbf{H-L7b:} exhaustive enumeration of all loss patterns of up to \texttt{W} = 8 consecutive lost ACKs and of reboot epochs (0 to 3 boots in the window): the base accepts the next alert exactly when the gap is at most \texttt{W}, and rejects (and never advances past) a gap above \texttt{W}. Cost: bytes (none added to the 24-byte provisional and 97-byte frames: the index is the existing counter) and hash operations per alert and per resynchronisation (ANALYTICAL).

\subsubsection*{15.4 L8: flood containment (bounded model and simulation)}

Base rule: per-vehicle token bucket of \texttt{R} = 6 high-confidence forwards per hour and 1 in any 60 s (spacing); plausibility: distance between consecutive reported positions at most 1.5 times (reported speed times elapsed time) plus 200 m; a repeated identical position within 10 minutes is implausible; a vehicle with two implausible or over-rate alerts in an hour enters quarantine, in which only its next alert in each 10 minutes is forwarded, as low confidence. \textbf{H-L8a (FORMAL, exhaustive bounded model):} over a horizon of one hour in 60 slots of 1 minute and every alert pattern of one compromised vehicle with up to 4 alerts per slot (explored by dynamic programming over the state of the bucket and the quarantine), the number of alerts forwarded as high-confidence is at most \texttt{R} = 6. \textbf{H-L8b (SIMULATED plus real rates):} legitimate alert streams with the alert rate per vehicle-hour taken from the VZCrash validation split (triggered events alarmed at the 1\% target, MEASURED\_SOFTWARE; 10 events per vehicle-hour as in the ns-3 scenarios), 2,000 simulated vehicle-hours, seed 42, positions consistent with a speed of 0 to 30 m/s: no legitimate alert is quarantined. Supplement unless H-L8a is proved.

\subsubsection*{15.5 L9: dispatch rule}

Descriptive, not in any family. Inputs: \texttt{r3.t3.a1.retraction\_share} with its interval and \texttt{r3.t3.a1.earlier\_ms\_median} with its interval (1\% target); the L2 result (not supported), so no corroboration is used. Options: act on the provisional alert, wait for the confirm, or pre-alert only (a pre-alert costs a fraction 0.1 of a dispatch; stated assumption). The break-even ratio \texttt{rho*} = (cost of a false dispatch) / (cost of one minute of delay) at which acting on the provisional alert and waiting have the same expected cost, with intervals by resampling the two inputs from their stored 95\% intervals (normal approximation, seed 42). No superiority claim.

\subsubsection*{15.6 L10: leakage audit}

The window-identity rule of C1 is applied unchanged in kind (exact hash of the float32 bytes of a window). Benchmark 1, UCI HAR: the published split (21 train and 9 test subjects, 128-sample windows with 50\% overlap, \texttt{Inertial Signals}): identical windows across splits and subject leakage (\texttt{subject\_train} against \texttt{subject\_test}); the headline of the benchmark is a classifier accuracy, so H-L10b is NOT\_RUN unless a leak is found (a random forest on the 561 features is then run as a stated baseline). Benchmark 2, SisFall (23 young adults, 15 older adults; files per subject, activity and trial; 200 Hz accelerometer): there is no published split, so the audit counts identical windows (1 s, 200 samples, step 100) within the dataset and between a random window split (80\% / 20\%, seed 42; the common protocol, stated); and trials that are byte-identical across subjects. \textbf{H-L10a:} at least one of the two benchmarks has more than 1\% of its test windows identical to a window of another split (random split for SisFall, published split for UCI HAR). If neither does, that is reported as a finding that narrows C1. \textbf{H-L10b:} where a baseline is run, removing the leaked windows changes its accuracy by more than 1 point; NOT\_RUN if no leakage is found.

\subsubsection*{15.7 L11: mirror-flip detection from turns (ADDED-AFTER-SESSION-1; own family)}

\textbf{Mechanism (rotation-invariant form).} For one event the up direction \texttt{z} is the unit mean acceleration; the horizontal dynamic acceleration \texttt{a\_h(t)} is the acceleration minus its mean, projected perpendicular to \texttt{z}; the yaw rate is \texttt{w(t) = gyro(t) . z} (rad/s); the speed \texttt{v(t)} (m/s) is the GPS speed interpolated to 100 Hz. In a right-handed frame \texttt{a\_h = (dv/dt) f + v w (z x f)} with \texttt{f} the forward direction. Two estimates of \texttt{f} are made in the sensor frame: \texttt{f\_L} from the longitudinal part (least squares of \texttt{a\_h} on \texttt{dv/dt} over samples with \texttt{|dv/dt|} above \texttt{A\_LONG}; \texttt{f\_L} is the unit vector of the coefficient vector) and \texttt{f\_T} from the turn (least squares of \texttt{a\_h} on \texttt{v w} over samples with \texttt{|w|} above \texttt{W\_MIN} and \texttt{v} above \texttt{V\_MIN}; \texttt{l\_T} is the unit vector of the coefficient vector, the lateral direction; \texttt{f\_T = l\_T x z}). The agreement of an event is \texttt{a\_e = f\_T . f\_L}, in [-1, 1]. It does not depend on how the sensor is rotated in the vehicle, and it changes sign under a mirror image of the frame. An event is \textbf{valid} if it has at least \texttt{N\_SEG} = 100 samples (1 s) of turn and of longitudinal change in the stated thresholds. \textbf{Decision:} an installation drives valid events until \texttt{N} = 20 valid events are collected (N is fixed here); the mean agreement \texttt{A} over them is computed; a mirror flip is declared if \texttt{A < THETA}. \textbf{Correction (H-L11b):} if a flip is declared, the lateral axis of the L1 rotation (its \texttt{y\_hat}) is negated before the side rule; for a single-axis mirror of x, y or z this restores the side decision (checked in the unit tests with noise-free signals). L1 itself is unchanged. The statistic does not need the L1 rotation for the VZCrash part, which has no calibration drive; this is a deviation from the task text ("after L1's rotation") and is stated in every table.

\textbf{Data (a), real signals with synthetic transformations (MEASURED\_SOFTWARE for the signals; transformations synthetic, stated).} VZCrash TRAIN split, normal-driving events (label 2). These events were used for fitting in earlier rounds and were never a held-out set for any hypothesis of rounds 3 to 5. Split by hash: \texttt{half\_l11(id) = int(sha256("l11|" + event\_id), 16) mod 2}; half 0 = DEVELOPMENT, half 1 = HELD-OUT. An installation is a group of consecutive valid events (in order of hashed event id) from its half; the groups are disjoint; up to \texttt{J} = 100 installations in the held-out half (fewer if the half has fewer valid events; the number is reported; if fewer than 40 installations can be formed, H-L11a is reported as not testable and not supported). Each group is evaluated under every transformation: identity; proper rotations (yaw 90, yaw 180, tilt 25 degrees about x, 180 degrees about x) applied to both channels (\textbf{unflipped installations}); mirror flips of x, y or z applied to both channels (\textbf{both-channel mirrors}); accelerometer-only flips of x, y, z; gyro-only flip of z. \textbf{Criterion transformations:} the three both-channel mirrors and the accelerometer-only flips of x and y. Expected and stated before the run: an accelerometer-only flip of z and gyro-only flips of x or y cannot be detected by this statistic (no horizontal signature); they are reported and are not in the criterion.

\textbf{H-L11a:} on the held-out half, with \texttt{N} = 20 valid events, the sensitivity is at least 0.95 for every criterion transformation and the false-detection rate on the unflipped installations (identity and the four rotations pooled) is at most 0.02. Verdict: point criteria hold and the largest exact one-sided binomial p of H0 "sensitivity at most 0.95" over the five criterion transformations is below 0.05 (own family of one; no Holm over the 13). False detection is reported with its exact 95\% upper limit. Descriptive: the share of events that are valid, and the number of clips and the driving time (16 s per clip) needed for 20 valid events.

\textbf{Development (at most 5 logged iterations, half 0):} items: \texttt{W\_MIN} (5, 10, 15 deg/s), \texttt{V\_MIN} (5, 10 m/s), \texttt{A\_LONG} (0.3, 0.6 m/s per second), \texttt{THETA} (0.0, 0.1, -0.1). Defaults: 10 deg/s, 5 m/s, 0.3 m/s per second, 0.0. Selection: the iteration with the highest development value of (minimum sensitivity over the criterion transformations) minus 10 times the false-detection rate; ties go to the iteration nearest the default.

\textbf{Data (b), simulated hardware-format logs (SIMULATED\_HW\_FORMAT, supplementary):} \texttt{r5\_l1.py} installations with mirror flips of x, y and z (seeds 70261117 development, 80261118 held-out), calibration drives extended with a turning drive (a left and a right turn of 0.15 g at 10 m/s, vehicle frame x forward, y left, z up, algebraically right-handed); \textbf{H-L11b:} the side accuracy after L1 plus L11 under single-axis mirror flips is at least 0.90 on the held-out installations (20 events per side), against about 0.5 with L1 alone. Simulated only: supplement (G0 rule 4).

\textbf{Code verification (G0 rule 3):} unit and property tests (an unflipped right-handed frame is never declared flipped by a single turn; flipping twice returns the original; the agreement is invariant under rotations), a known-answer test (noise-free synthetic turn), at least 20 mutants with at least 80\% killed, and a fresh-context review of the new modules before the freeze. \textbf{Literature row before the run;} if earlier work detects handedness errors in this way, L11 is reported as an application, not a contribution. Paper: if supported, in the main text next to L1 as "added after session 1, own family", with the L1 mirror-flip limitation kept; if not, the negatives register.

\subsubsection*{15.8 Cut order and freeze}

Cut order for session 2: L9, L8b, L11b. Never cut S2, L6, L7a, L10a, L11a, R, D. Each algorithm: development (at most 5 logged iterations where it has tunable items), freeze by hash with \texttt{r5\_common.freeze\_entry}, held-out run once (\texttt{assert\_fresh}). The hashes of the L11 halves are written to \texttt{results/tables/r5\_split\_hashes\_l11.json} before the held-out run; the L6 year split is fixed by the years above.


\subsubsection*{15.9 Addendum S2.1 (L11; before any L11 run, including the build of any cache)}

\begin{itemize}
\item \textbf{Reference sign.} The native frames of the VZCrash devices need not follow the sign convention of the relation in 15.7 (a device may report the gyro with the opposite sign). The reference sign \texttt{s0} in \{+1, -1\} is the sign of the median agreement of the identity-transformed DEVELOPMENT installations (development half only); it is frozen with the other parameters and every agreement is multiplied by \texttt{s0} before the decision. If the absolute value of that median is below 0.3 on development, the dataset is declared inconsistent with the mechanism and H-L11a is reported as not supported.
\item \textbf{Event cache.} To keep the arrays in memory, only events that are valid under the loosest development thresholds (\texttt{W\_MIN} 5 deg/s, \texttt{V\_MIN} 5 m/s, \texttt{A\_LONG} 0.3 m/s per second, 100 samples each) under the identity transformation are cached. Validity under stricter thresholds is a subset; validity under the transformations is evaluated on the cached events (for rotations and both-channel mirrors the thresholds act on invariant quantities).
\item \textbf{Installation order.} Installations are formed from the valid events of a half in ascending order of \texttt{sha256("l11-order|" + event\_id)}, groups of consecutive valid events; installations are disjoint; the number is \texttt{min(100, floor(valid / 20))}.
\item \textbf{L11b} (simulated) is cut if time is short and then reported as NOT\_RUN.
\end{itemize}


\subsubsection*{15.10 Addendum S2.2 (L11; before any L11 run)}

While the unit tests were written (synthetic drives only; no VZCrash event had been read) it was found that the unit mean acceleration of the whole event, named in 15.7 as the up direction \texttt{z}, is tilted by the mean lateral acceleration of the turn (about 10 degrees for a 6 s turn of 0.2 g) and spoils both estimates. The up direction is therefore the unit mean acceleration of the quasi-steady samples: samples with a gyro norm below 3 deg/s and \texttt{|dv/dt|} below \texttt{A\_LONG}; if fewer than 50 such samples exist, all samples are used. Nothing else in 15.7 changes.

